% Generated by maint/texbuild/theory_tex.py from theory/workbook/vertical_time_compression.md.
% Edits here are overwritten. Edit the markdown, then run the script again.
\chapter{\texorpdfstring{Vertical time compression: floors stacked into one sweep}{Vertical time compression: floors stacked into one sweep}}
\label{chap:vertical-time-compression}

\textbf{Purpose:} the algebra of the record machine's stacked floors: what a stack is, what it compresses, what it leaves as it is, and what is proved and measured about it.

\textbf{Scope:} the record machine (M10 and A13 in \hyperref[chap:engine-table]{engine\_\allowbreak{}table.md}): \texttt{keymath} imprints a program, \texttt{key\_\allowbreak{}schedule} lays it out, \texttt{cycle} sweeps it. The shape is Doug's (24 September): a floor, the operations in order above it, the next floor, and so on up, with the top conjunction performing the whole transform.

\section{\texorpdfstring{The objects}{The objects}}

\begin{itemize}
\item \textbf{Lane.} A lane ℓ reads its atom a\_\allowbreak{}ℓ, the fields of its records.
\item \textbf{Floor.} A floor S\_\allowbreak{}k is the set of registers live at a cut between two rounds. It is the state one round leaves and the next reads.
\item \textbf{Round.} A round f\_\allowbreak{}k is a straight-line list of steps. It reads S\_\allowbreak{}\{k−1\} and leaves S\_\allowbreak{}k. Its steps fill the register file while they run. At the floor, register reuse has freed everything except S\_\allowbreak{}k.
\item \textbf{Stack.} The stack is F = f\_\allowbreak{}n ∘ … ∘ f\_\allowbreak{}1, imprinted once as one key. The imprint derives every register's width across the whole composition before any lane runs.
\item \textbf{Sweep.} A sweep applies F to every lane in one launch: R\_\allowbreak{}ℓ = F(a\_\allowbreak{}ℓ).
\end{itemize}

\section{\texorpdfstring{The algebra}{The algebra}}

\begin{enumerate}
\item \textbf{Composition regroups and keeps its order.}

\begin{itemize}
\item (h ∘ g) ∘ f = h ∘ (g ∘ f), but g ∘ f ≠ f ∘ g in general.
\item A stack can be cut into programs at any floor and the programs chained, and every result stays the same.
\item \textbf{Proved} for tables (M10): a table read through a table equals the ops, and the reversed order differs on some lane.
\item \textbf{Proved} for floors (below): 700 floors run as one program equal the same floors run as 700 chained programs, record for record, at four lane counts.
\end{itemize}
\item \textbf{The file is bounded by the widest floor, not by the step count.}

\begin{itemize}
\item A register is live from the step that writes it to its last reader. With reuse on, the file needs the most limbs live at any one step. That is at most the widest floor plus the round in flight.
\item The step count has no bound of its own. \texttt{ENGINE\_\allowbreak{}RECORD\_\allowbreak{}STEPS\_\allowbreak{}MAX} held one sign per step. Since 24 September each sign sits beside its register. The file (\texttt{ENGINE\_\allowbreak{}RECORD\_\allowbreak{}LIMBS\_\allowbreak{}MOST}, 256 limbs) is now the only bound.
\item \textbf{Measured:} 4,204 steps in an 8-limb file (the stack test), and in a 12-limb file with its words carried as 33-bit fields (the chained measure).
\end{itemize}
\item \textbf{Widths are derived, never declared.}

\begin{itemize}
\item Each step's width follows from its operands' widths.
\item A register the imprint can show is never negative carries that fact forward:

\begin{itemize}
\item a field read unsigned, a constant, an absolute value, a gcd, a table entry;
\item a sum, product, quotient or xor of two such registers;
\item an and with one such register.
\end{itemize}
\item With that fact, an and is no wider than its never-negative operand, and an xor of two never-negative registers is no wider than the wider one. A stack of rounds on fixed-width words then keeps its words at their width through every floor, instead of growing a bit per round.
\item Every width is checked against the values: the bitwise oracle rebuilds each output bit from the raw fields and requires the sign to hold 64 bits past the written width.
\end{itemize}
\item \textbf{Time: what the stack compresses.}

\begin{itemize}
\item \textbf{Chained}, each of the n floors pays a launch, a synchronization, and a round trip of its state through device memory, on top of its steps.
\item \textbf{Stacked}, the lane reads its atom once, runs every step in registers, and writes its record once. The time between floors is gone.
\item The lanes run side by side on the device. Until the device is full, a stack over L lanes costs one tower's depth for the whole batch.
\item \textbf{Measured} (below): stacked is 13 to 22 times faster than chained.
\end{itemize}
\item \textbf{Time: what the stack leaves as it is.}

\begin{itemize}
\item Each lane still evaluates every step in order. Once the device is full, a sweep's time grows linearly with the number of steps.
\item The only ways to shorten that are fewer steps:

\begin{itemize}
\item a table (a sub-function of up to 32 input bits becomes one lookup, and composition keeps its order);
\item narrower registers (fewer limbs a step).
\end{itemize}
\end{itemize}
\item \textbf{Exactness.}

\begin{itemize}
\item Every step is exact integer arithmetic.
\item The host evaluates the same key with the exact integer library.
\item A device that equals the host word for word proves the port, for a stack as for any program.
\end{itemize}
\end{enumerate}

\section{\texorpdfstring{Proved and measured}{Proved and measured}}

\begin{itemize}
\item \textbf{The stack test} (\texttt{test/\allowbreak{}engine/\allowbreak{}record\_\allowbreak{}bitwise\_\allowbreak{}test}). 700 floors of one round each, over four 32-bit words: xor, and, sum, xor with the floor's constant, and a 32-bit wrap.

\begin{itemize}
\item 4,204 steps, one program, register reuse on, an 8-limb file.
\item On 1,024 lanes the device equals the host word for word. Both equal the same rounds run on the CPU's own 64-bit two's complement at every tapped floor.
\end{itemize}
\item \textbf{Stacked against chained.} The same 700 floors, run as one program of 4,204 steps (a 12-limb file) and as 700 programs of 13 steps. Each chained sweep reads the previous sweep's records from device memory. RTX 3070, the second sweep of each timed:

{\footnotesize
\setlength{\tabcolsep}{6pt}
\begin{longtable}{>{\RaggedRight\arraybackslash}p{\dimexpr0.2032\linewidth-2\tabcolsep\relax}>{\RaggedRight\arraybackslash}p{\dimexpr0.2202\linewidth-2\tabcolsep\relax}>{\RaggedRight\arraybackslash}p{\dimexpr0.2392\linewidth-2\tabcolsep\relax}>{\RaggedRight\arraybackslash}p{\dimexpr0.1583\linewidth-2\tabcolsep\relax}>{\RaggedRight\arraybackslash}p{\dimexpr0.1792\linewidth-2\tabcolsep\relax}}
\toprule
lanes & stacked, one sweep & chained, 700 sweeps & ratio & records \\
\midrule
\endhead
1,024 & 4,274 µs & 55,899 µs & 13.1× & equal \\
16,384 & 4,122 µs & 71,419 µs & 17.3× & equal \\
262,144 & 26,048 µs & 572,534 µs & 22.0× & equal \\
1,048,576 & 98,783 µs & 2,088,411 µs & 21.1× & equal \\
\bottomrule
\end{longtable}
}

\begin{itemize}
\item At 1,024 and 16,384 lanes the stacked sweep takes the same time: the device is not yet full, and the batch costs one tower's depth.
\item From 262,144 lanes the time grows with the lanes, and the ratio settles near 21×. That ratio is the cost of the time between floors, which the stack removes.
\item At 1,048,576 lanes the stacked sweep evaluates 4.4 × 10\textasciicircum{}9 steps in 98.8 ms, about 4.5 × 10\textasciicircum{}10 steps a second.
\item The measure's source is held outside the repository and has not been rerun from it.
\end{itemize}
\end{itemize}

\section{\texorpdfstring{What the stack is not}{What the stack is not}}

\begin{itemize}
\item It does not evaluate fewer steps than the program has. The depth stays; the compression is of the time between floors and across the lanes.
\item It does not change a result: a stack and its chained programs give the same records.
\end{itemize}

\section{\texorpdfstring{The two towers}{The two towers}}

Doug's (24 September): a second tower stacked over the first one's boundary, inverted, and the two collapse together. One is the crystal, the tower of the data's lifted floors. The other is the tower of operations, the stack above.

\begin{itemize}
\item \textbf{The crystal's tower T} is the 5/3 integer lifting of \texttt{src/\allowbreak{}engine/\allowbreak{}base/\allowbreak{}tower/\allowbreak{}tower.\allowbreak{}cu} (A14 in \hyperref[chap:engine-table]{engine\_\allowbreak{}table.md}). One level over a line of samples x:

\begin{itemize}
\item the high d\_\allowbreak{}j = x\_\allowbreak{}\{2j+1\} − ⌊(x\_\allowbreak{}\{2j\} + x\_\allowbreak{}\{2j+2\}) / 2⌋;
\item the low s\_\allowbreak{}i = x\_\allowbreak{}\{2i\} + ⌊(d\_\allowbreak{}\{i−1\} + d\_\allowbreak{}i + 2) / 4⌋;
\item an edge repeats its neighbor.
\end{itemize}
\item \textbf{Its inverse T⁻¹} runs the same two lines backward and in the other order: x\_\allowbreak{}\{2i\} = s\_\allowbreak{}i − ⌊(d\_\allowbreak{}\{i−1\} + d\_\allowbreak{}i + 2) / 4⌋, then x\_\allowbreak{}\{2j+1\} = d\_\allowbreak{}j + ⌊(x\_\allowbreak{}\{2j\} + x\_\allowbreak{}\{2j+2\}) / 2⌋.
\item \textbf{Floor division is a record floor.} For every integer v and k ≥ 0:

⌊v / 2\textasciicircum{}k⌋ = EXACT\_\allowbreak{}QUOTIENT(v − AND(v, 2\textasciicircum{}k − 1), 2\textasciicircum{}k)

\begin{itemize}
\item The and reads v's two's complement without end. With the never-negative mask 2\textasciicircum{}k − 1 it gives v's residue modulo 2\textasciicircum{}k, in [0, 2\textasciicircum{}k), for a negative v as for a positive one.
\item v less its residue is a multiple of 2\textasciicircum{}k. The exact quotient divides it with nothing left over, and the quotient rounds toward −∞, as \texttt{tower\_\allowbreak{}floor\_\allowbreak{}shift} does.
\item The residue is no wider than the mask, k bits, by the never-negative width rule (3).
\end{itemize}
\item \textbf{T and T⁻¹ are stacks of record floors.} Every step of a lifting level is a sum, a difference, a constant or that floor division. A level is a floor of the record machine, L levels are L floors, and the inverse is L more.
\item \textbf{The operation tower over the crystal.} A program F that reads the crystal's lifted floors, and the inverse T⁻¹ that brings them back, compose into one stack F ∘ T⁻¹ by the regrouping law (1). No step between them leaves the lane, and T⁻¹ ∘ T is the identity on the machine exactly.
\item \textbf{Proved} (\texttt{test/\allowbreak{}engine/\allowbreak{}record\_\allowbreak{}bitwise\_\allowbreak{}test}): one 5/3 level over 8 signed 16-bit samples and its inverse, 129 steps as one program, register reuse on, a 12-limb file, 4,096 lanes of edge-shaped samples.

\begin{itemize}
\item The device equals the host word for word.
\item The forward floor's 4 lows and 4 highs equal tower.cu's formulas, computed on the CPU, on every lane.
\item The inverse floor returns all 8 samples exactly on every lane.
\end{itemize}
\end{itemize}

\section{\texorpdfstring{Open}{Open}}

\begin{itemize}
\item \textbf{A lane's own index as a register.} A lane now reads its inputs only from records. An operation giving the lane's number as a value would let one shared atom stand for a whole range of inputs, with the lanes enumerating the range and no input stored per lane.
\item \textbf{The latch.} A device reduction that returns the first lane whose output meets a condition. Only that lane's index comes back to the host.
\item \textbf{Depth by tables.} Which rounds' sub-functions fit a 32-bit index, and how much depth that removes.
\item \textbf{The chained cost apart.} Launch, synchronization and the state's round trip through memory, each measured on its own.
\item \textbf{The neighbor gather.} A lifting level reads each sample's neighbors, and L levels in one lane read the cone of samples under that lane, widening with every level. A lane reads at most 3 members through the index, one record from each. Whether the cone is laid into a lane's record, or the levels split across sweeps where the cone outgrows it, is open.
\item \textbf{What F ∘ T⁻¹ saves.} Whether a linear F passes through T⁻¹ or folds into it, in fewer steps than the two apart, is not measured.
\item \textbf{The whole crystal as one stack.} The proof above is one level along one line. All levels along all four axes as record floors, against tower.cu's own crystal, is not built.
\end{itemize}
